% ---------------------------------------------------------------------------
% Chassis / CAD part (3D design, written by the 3D-design team member).
% Moved out of diplomarbeit_chap_3.tex on 26.11.2026 so that the hardware text
% could be rewritten in English. Content unchanged.
% TODO (3D design): translate to English and update after the first print:
%  - PCB standoffs were raised by 4 mm (top at Z = 29.5 mm instead of 25.5 mm),
%    because the PCB underside (D14, test points, header pins) would touch the battery;
%  - the part was printed in PETG on a Bambu Lab P1S WITH build-plate supports
%    (first attempt without supports failed in the first layer, see Section~\ref{sec:hw_print}).
% ---------------------------------------------------------------------------
\begin{otherlanguage}{ngerman}
\sisetup{locale=DE}
\todo[inline]{3D-Design: Abschnitt auf Englisch übersetzen und an den gedruckten Stand anpassen (Standoffs Z~=~\SI{29.5}{\milli\metre}, Druck mit Stützen) -- siehe Kommentar in \texttt{diplomarbeit\_chap\_3\_cad.tex}.}
\section{CAD-Konstruktion und Gehäuse-Evolution (Revision v1 bis v9.0)}
\label{sec:cad_konstruktion}

Die mechanische Konstruktion des Fahrzeugchassis wurde in Autodesk Fusion~360 sowie automatisiert über Python (\texttt{build123d}) entwickelt und über neun Revisionen hinweg schrittweise verfeinert. Frühe Entwürfe (v1/v2) wiesen noch zu geringe Wandstärken, blockierte Schraubenbohrungen am Liniensensor sowie abgeschattete Sichtkegel der hinteren Sensoren durch Räder und Motoren auf.

\begin{figure}[htbp]
  \centering
  \includegraphics[width=0.96\textwidth]{v3_LiDar_V2_engineering_views.png}
  \caption{Ingenieuransichten der Sensor-Pod-Geometrie mit hinter den Antriebsrädern positionierten Heck-Sensoren.}
  \label{fig:cad_views_v3}
\end{figure}

In der finalen Ausbaustufe (**Revision~v9.0 Final DMU**) wurde das gesamte Fahrzeug als einteiliges, stützstrukturfrei druckbares Unibody-Chassis (\SI{0.0}{\gram} Support, Volumen \SI{218.8}{\centi\metre\cubed}, Masse ca.~\SI{271}{\gram} in PLA/PETG) umgesetzt. Abbildung~\ref{fig:v9_master_dmu} zeigt das vollständige digitale Versuchsmodell (\gls{dmu}) inklusive aller 121~Einzelkörper.

\begin{figure}[htbp]
  \centering
  \includegraphics[width=0.98\textwidth]{v9_Final_DMU_Master_Preview.png}
  \caption{Revision~v9.0 Final DMU: Isometrische Gesamtansicht, Draufsicht, Seitenansicht und Unteransicht inklusive 8x-LiDAR-Sichtkegeln (\SI{25}{\degree}), 24-mm-Steckerdummies (rot) und 3D-Kabelbaum (orange).}
  \label{fig:v9_master_dmu}
\end{figure}

\subsection{8x-LiDAR-Sensor-Pods und Schraubendome}
Um alle acht \emph{VL53L0/1XV2}-Platinen (\qtyproduct{25.0 x 10.5 x 1.6}{\milli\metre}, Lochabstand \SI{19.5}{\milli\metre}) stoßgeschützt aufzunehmen, sitzt jedes Modul um \SI{1.90}{\milli\metre} tief in einer eigenen Sensortasche unter einem durchgehenden \SI{3.85}{\milli\metre} bis \SI{4.30}{\milli\metre} starken Massivdach. Die Befestigung erfolgt über jeweils zwei massive Schraubendome (\SI{6.5}{\milli\metre} Außendurchmesser) mit \SI{6.0}{\milli\metre} tiefen M2-Sacklöchern ($\varnothing\,\SI{2.0}{\milli\metre}$). Hinter jedem Sackloch verbleiben \SI{3.50}{\milli\metre} geschlossenes Vollmaterial, sodass M2x6-Schrauben zuverlässig greifen, ohne die Rückwand zu durchstoßen.

Für die sechspoligen DuPont-Stecker auf der Sensorrückseite stehen geschützte Steckertunnel mit \SI{26.0}{\milli\metre} Tiefe zur Verfügung. Alle sechs seitlichen bzw. schrägen Pods reichen mit \SI{15.5}{\milli\metre} Einbindetiefe lückenlos in die \SI{3.5}{\milli\metre} starke Hauptwannenwand, während die Kabelkanäle durch den \SI{4.0}{\milli\metre} hohen geschlossenen Wandsteg darüber die Struktursteifigkeit erhalten.

\subsection{Batteriefach-Toleranzen, Hauptplatinen-Standoffs und FDM-Regeln}
Die Hauptplatine ruht auf vier konisch verstärkten Standoffs ($\varnothing\,\SI{10.6}{\milli\metre}$, Auflageebene $Z = \SI{25.5}{\milli\metre}$) mit \SI{6.0}{\milli\metre} tiefen Sitzen für M3-Messing-Gewindeeinsätze ($\varnothing\,\SI{4.0}{\milli\metre}$, radiale Restwandstärke \SI{3.30}{\milli\metre}). Eine konstruktive Schutzoperation stellt sicher, dass kein Kabelkanal die Standoff-Säulen seitlich anschneidet.

Das rechts in der Wanne liegende Akkufach für das 2S-Li-Ion-Pack (\qtyproduct{72.0 x 39.5 x 19.5}{\milli\metre}) bietet einen lichten Bauraum von \qtyproduct{78.0 x 43.5 x 21.5}{\milli\metre}. Daraus ergeben sich reale Einbautoleranzen von \SI{+2.0}{\milli\metre} seitlich und vertikal sowie \SI{+3.0}{\milli\metre} in Längsrichtung, sodass der Akku werkzeuglos ohne Klemmen entnommen werden kann. Eine \SI{3.2}{\milli\metre} starke Trennrippe schützt den Akku mechanisch vor dem \SI{24.0}{\milli\metre} langen Stecker- und Biegeraum des Liniensensors (Gesamtabstand \SI{8.0}{\milli\metre}). Für den USB-C-Stecker an \texttt{J4} bleibt oberhalb des Front-Pod-Dachs ($Z = \SI{23.0}{\milli\metre}$ bis \SI{33.0}{\milli\metre}) ein \SI{16.0}{\milli\metre} breiter Einsteckkorridor dauerhaft frei.

\subsection{Rechnerische Crash-Abschätzung}
Zur Bewertung der Aufprallsicherheit wurde eine analytische Stoßabschätzung für einen ungebremsten Frontalaufprall mit der Fahrzeuggesamtmasse $m = \SI{1.20}{\kilogram}$ und einer elastisch-plastischen Pufferzone $s = \SI{2.0}{\milli\metre}$ am \SI{6.0}{\milli\metre} starken Front-Bumper durchgeführt. Über den Energieerhaltungssatz ergibt sich die mittlere Stoßkraft zu:
\begin{equation}
  F_{\text{Crash}} = \frac{\frac{1}{2}\,m\,v^2}{s}
  \label{eq:crash_kraft}
\end{equation}
Bei einer typischen Einpark- bzw. Suchgeschwindigkeit von $v_1 = \SI{5.0}{\kilo\metre\per\hour}$ (\SI{1.39}{\metre\per\second}) resultiert $F_1 \approx \SI{579}{\newton}$; selbst im Extremfall bei $v_2 = \SI{9.0}{\kilo\metre\per\hour}$ (\SI{2.50}{\metre\per\second}) beträgt die Stoßkraft $F_2 = \SI{1875}{\newton}$. Bezogen auf die tragende Prallquerschnittsfläche von $A = \SI{612}{\milli\metre\squared}$ und eine zulässige Biege-/Druckspannung von $\sigma_{\text{zul}} = \SI{35}{\mega\pascal}$ für PLA/PETG ergeben sich Sicherheitsfaktoren von $\text{SF}_1 = 37{,}0$ bzw. $\text{SF}_2 = 11{,}4$, die deutlich über dem Zielwert $\text{SF} \ge 3{,}0$ liegen.

\end{otherlanguage}
